\chapter{Bilan : le socle, ses limites, sa reprise}
\label{ch:bilan}

\questionchapitre{Que reste-t-il une fois le contrat achevé, dans quel état, et à quelles
conditions ces développements peuvent-ils servir au-delà du projet qui les a financés ?}

\section{Récapitulatif des réalisations}

\begin{table}[H]
\centering
\caption{Les réalisations de la période, par domaine.}
\label{tab:realisations}
\begin{tabularx}{\textwidth}{@{}L{3.4cm}Y@{}}
\toprule
\textbf{Domaine} & \textbf{Réalisation} \\
\midrule
Données & Chaîne d'ingestion automatisée de quatre jeux de données issus de trois fournisseurs
publics ; entrepôt en architecture médaillon sur base relationnelle avec extensions temporelle
et géospatiale ; rattachement spatial des stations à leur maille climatique ; enrichissement par
le référentiel hydrogéologique ; environ \num{658} millions de lignes en couche brute \\
Fiabilité & Idempotence des traitements, traçabilité des lignes écartées, contraintes de
concurrence par clés, reprise après interruption, sauvegarde et restauration vérifiées sur un
entrepôt représentatif \\
Indicateurs & Indices standardisés de niveau, de débit, de précipitation, de température et de
bilan hydrique ; production à l'échelle nationale sur quatre fenêtres ; protocole de validation
à critère pré-établi ; deux décisions méthodologiques établies par la mesure \\
Exploration & Observatoire spatial des nappes et des rivières ; module climatique national avec
historique depuis 1950, épisodes de sécheresse et export \\
Modélisation métier & Laboratoire de modèles à fonction de transfert ; calibration automatique
guidée par la nature de l'aquifère ; évaluation selon quatre critères normalisés ; diagnostics
préalables et postérieurs ; scénarios prospectifs à bornes physiquement dérivées \\
Apprentissage & Catalogue d'une douzaine d'architectures de prévision profonde ; optimisation
d'hyperparamètres ; traçabilité des expériences ; suivi d'entraînement en temps réel ;
configurations prêtes à l'emploi par cas d'usage métier \\
Explicabilité & Attribution multi-méthodes ; contrefactuels sous contrainte physique avec
couplage thermodynamique et validation par modèle indépendant ; détection non supervisée
d'influence anthropique \\
Ingénierie & Deux dépôts conteneurisés, reproductibles et documentés pour la reprise ;
environnements de développement et de production séparés \\
\bottomrule
\end{tabularx}
\end{table}

\section{Bilan au regard de la mission initiale}

La mission portait sur l'explicabilité des modèles de prévision appliqués aux eaux souterraines.
Elle a été instrumentée de bout en bout ; elle n'a pas été conclue.

\textbf{Ce qui a été atteint} : l'outillage complet que le sujet supposait, jusqu'aux méthodes
d'explication elles-mêmes. Parmi elles, les contrefactuels sous contrainte physique s'attaquent
directement à ce qui rend l'explicabilité délicate dans ce domaine, à savoir qu'une explication
physiquement incohérente n'est pas une explication. Ces développements restent préliminaires :
ils sont opérationnels et prêts à être éprouvés, ils ne sont pas validés.

\textbf{Ce qui n'a pas été atteint} : l'évaluation quantitative comparée des méthodes, faute de
temps une fois l'infrastructure construite. C'est le travail le plus immédiatement accessible à
qui reprendra ces développements, parce que tout ce qu'il suppose existe désormais.

\section{Ce qui rend ce socle réutilisable}

Affirmer qu'un développement constitue un « socle » n'a de valeur que si l'on peut désigner les
propriétés qui le rendent réutilisable et montrer qu'elles sont acquises. Elles sont au nombre
de quatre.

La couche analytique de l'entrepôt se consomme en SQL, et rien d'autre. Un utilisateur n'a besoin
de connaître ni l'outil d'ingestion, ni la bibliothèque de transformation, ni l'orchestrateur :
un carnet de calcul, un tableur connecté ou une application tierce s'y branchent également. La
conséquence est structurante : \textbf{l'amont peut être remplacé sans que l'aval s'en
aperçoive}. Un outil d'ingestion devenu obsolète se remplace sans toucher aux consommateurs.
C'est cette propriété qui distingue une infrastructure d'un script.

Le cœur de calcul de la plateforme, lui, ne dépend ni du service web ni de l'interface. La contrainte a été posée
dès l'origine, au chapitre~\ref{ch:observatoire}, et son effet sur la valeur résiduelle des
travaux est direct : \textbf{même si l'application web était abandonnée, les méthodes
resteraient utilisables}. La contribution n'est pas prisonnière de son emballage.

Les indicateurs sont calculés une fois pour tous les consommateurs, sans second calcul
susceptible de diverger, et un test de contrat échoue si les deux dépôts cessent de s'accorder
(chapitre~\ref{ch:indicateurs}), sous la réserve d'exécution rappelée plus bas. La pile, enfin, est conteneurisée, ses volumes de
données survivent à l'arrêt des conteneurs, sa restauration a été vérifiée sur un entrepôt
représentatif et non sur une table jouet (chapitre~\ref{ch:entrepot}), et sa documentation a été
refondue en fin de période pour la reprise, jusqu'aux pièges qui produisent des résultats faux
sans émettre d'erreur.

Au-delà des deux dépôts, ces travaux instancient un enchaînement (figure~\ref{fig:patron}) qui
n'a rien de spécifique aux eaux souterraines.

\begin{figure}[H]
\centering
\begin{tikzpicture}[
  font=\sffamily\small,
  et/.style={draw=black!70, fill=black!3, minimum width=2.6cm, minimum height=1.3cm,
             align=center, font=\sffamily\scriptsize},
  fl/.style={-{Stealth[length=2.4mm]}, draw=black!70, thick},
  note/.style={font=\sffamily\scriptsize\itshape, color=black!55, align=center,
               text width=2.55cm},
]
\node[et] (e1) {\textbf{Sources}\\ouvertes,\\hétérogènes};
\node[et, right=0.7cm of e1] (e2) {\textbf{Entrepôt}\\consolidé,\\interrogeable};
\node[et, right=0.7cm of e2] (e3) {\textbf{Indicateurs}\\normalisés\\et validés};
\node[et, right=0.7cm of e3] (e4) {\textbf{Modèles}\\métier et\\appris};
\node[et, right=0.7cm of e4, very thick, draw=black] (e5) {\textbf{Explication}\\et\\scénarisation};

\foreach \a/\b in {e1/e2, e2/e3, e3/e4, e4/e5} \draw[fl] (\a) -- (\b);

\node[note, below=0.32cm of e2] {coût payé une seule fois};
\node[note, below=0.32cm of e3] {comparabilité garantie ici};
\node[note, below=0.32cm of e5] {condition de l'usage en décision};
\end{tikzpicture}
\caption{L'enchaînement mis en œuvre. Chaque étape suppose la précédente ; escamoter l'une
d'elles produit des résultats incomparables ou inutilisables.}
\label{fig:patron}
\end{figure}

L'enseignement principal de l'année tient dans la lecture de gauche à droite de cette figure :
\textbf{on ne peut pas commencer par la droite}. Investir dans le modèle avant de disposer d'une
donnée exploitable conduit à découvrir ensuite qu'elle ne permet pas de conclure.

\section{Ce qui est transmis aux travaux à venir}

Les travaux d'explicabilité se poursuivront au-delà de cette période. Le
tableau~\ref{tab:transmis} confronte ce qu'ils supposent à ce qui leur est effectivement
transmis.

\begin{table}[H]
\centering
\caption{Ce que supposent les travaux d'explicabilité à venir, et ce dont ils
disposent.}
\label{tab:transmis}
\begin{tabularx}{\textwidth}{@{}L{4.6cm}Y@{}}
\toprule
\textbf{Besoin} & \textbf{Disponible} \\
\midrule
Un jeu de données consolidé, joint au forçage climatique & Entrepôt national, partitions depuis
1967 pour les stations et 1950 pour le climat, rattachement spatial matérialisé \\
Des modèles à expliquer & Douze architectures de prévision profonde sous interface unique, avec
traçabilité des expériences \\
Une référence interprétable à laquelle confronter les explications & Modèles à fonction de
transfert calibrés, dont les paramètres ont un sens hydrogéologique \\
Des méthodes d'attribution & Cinq familles complémentaires, confrontables entre elles \\
Des explications actionnables & Contrefactuels sous contrainte physique, mettant en œuvre les
critères formulés par les hydrogéologues \\
Un moyen de valider une explication & Rejeu du scénario dans un modèle métier indépendant \\
Des cas d'étude qualifiés & Détection non supervisée d'influence anthropique, permettant
d'écarter les chroniques impropres \\
\bottomrule
\end{tabularx}
\end{table}

Ce qui n'est pas transmis tient en une phrase : le socle rend ces travaux possibles, il ne les a
pas menés. Le détail en est donné plus bas, aux points 3 et suivants.

\section{Limites et chantiers à reprendre}

Un successeur doit savoir ce que le socle ne fait pas, et ce qu'il faut engager pour qu'il le
fasse. Les deux questions n'en font qu'une : une limite qui se corrige est un chantier. La liste
suivante est ordonnée par priorité, la première ligne étant la seule qui conditionne la validité
de résultats à venir.

\begin{enumerate}
  \item \textbf{Les modèles métier sont calibrés sur une évapotranspiration inadaptée.} La
  chaîne des stations transporte encore la variable brute, environ deux fois supérieure à
  l'évapotranspiration de référence employée par la chaîne climatique. La mesure présentée au
  chapitre~\ref{ch:indicateurs} montre que le modèle n'absorbe pas cet écart : le paramètre qui
  pondère l'évaporation sature à sa borne, et l'ajustement en dépend. Les prévisions restent
  exploitables, mais les paramètres calibrés n'ont pas d'interprétation physique directe.
  Aligner la chaîne des stations est le premier chantier, et il précède toute campagne de
  calibration.

  \item \textbf{Aucun des deux dépôts n'exécute ses tests automatiquement.} Leurs chaînes
  d'intégration ne comportent qu'une étape de construction et une étape de déploiement. Sur la
  plateforme, cinq des \num{552} tests ne passent pas, quatre échecs et un test ignoré, pour des
  raisons documentées et sans rapport avec un défaut applicatif, mais la purge des journaux se
  trouve de fait non testée. Surtout, le contrat de classification décrit au
  chapitre~\ref{ch:indicateurs} n'est gardé par aucune exécution : les deux tables de référence
  sont bien identiques et le test qui les compare existe de part et d'autre, mais rien ne le
  lance. Tant que les vérifications dépendent d'une initiative humaine, la confiance dans les
  évolutions décroît avec le temps. C'est le chantier le moins coûteux et le plus rentable.

  \item \textbf{Aucun résultat de performance comparée n'est établi.} Les campagnes
  d'entraînement ont validé le dispositif, non classé les méthodes ; les jeux préparés et les
  points de sauvegarde ont été purgés lors de la passation. La campagne d'évaluation
  comparative, qui opposerait modèles métier et modèles appris sur un échantillon représentatif
  de stations et de familles d'aquifères, reste donc à conduire. L'outillage est prêt de part et
  d'autre, et le socle sera d'autant plus crédible qu'il aura servi à établir un résultat.

  \item \textbf{La couverture climatique s'arrête à la métropole.} La grille ingérée couvre la
  France métropolitaine et la Corse. Les stations ultramarines sont présentes dans l'entrepôt
  mais sans forçage climatique, ce qui les rend inutilisables pour la modélisation.

  \item \textbf{Certaines températures agrégées au niveau des stations} sont des extrêmes de
  moyennes journalières et non de véritables extrêmes journaliers ; l'exposition de ces derniers
  supposerait de modifier une table partitionnée.

  \item \textbf{La source climatique retenue est un compromis.} ERA5-Land a été préféré à une
  réanalyse plus fine mais non redistribuable, au bénéfice de la reproductibilité et du partage.

  \item \textbf{Le périmètre couvre la quantité, non la qualité} des eaux, également visée par
  le programme.
\end{enumerate}

Trois autres chantiers ne corrigent aucune limite du livrable, mais conditionnent son usage
au-delà du cercle actuel. Aucun n'est un verrou de recherche ; tous sont des décisions
d'investissement.

\begin{itemize}
  \item \textbf{Un accès élargi et gouverné.} L'accès nominatif convient à un outil de projet,
  non à un environnement partagé. Une politique d'accès, précisant qui accède, pour quels usages
  et avec quelle responsabilité sur les données, conditionne tout élargissement.

  \item \textbf{Un catalogue de données.} L'entrepôt est documenté pour qui sait où chercher.
  Un environnement d'outils suppose un catalogue décrivant les jeux disponibles, leur origine,
  leur fraîcheur et leurs conditions d'usage.

  \item \textbf{Le calcul mutualisé.} L'entraînement s'appuie aujourd'hui sur un serveur unique.
  Un usage partagé suppose une file d'attente et un partage de ressources.
\end{itemize}

\section{Extensions, par coût croissant}

\begin{table}[H]
\centering
\caption{Ce que coûterait chaque extension du socle, et ce qu'elle suppose.}
\label{tab:extensions}
\begin{tabularx}{\textwidth}{@{}lXl@{}}
\toprule
\textbf{Extension} & \textbf{Ce qu'elle suppose} & \textbf{Coût} \\
\midrule
Nouvelles stations, autre territoire & Rien : la couverture est déjà nationale et la chaîne est
paramétrée par configuration & nul \\
Nouveaux indicateurs standardisés & Un ajustement statistique et sa validation, selon le
protocole existant & faible \\
Nouvelles architectures de prévision & Une entrée au catalogue ; la fabrique de modèles et la
traçabilité sont génériques & faible \\
Qualité des eaux souterraines & De nouvelles sources et de nouveaux modèles de transformation ;
l'architecture est inchangée & moyen \\
Eaux de surface & La chaîne hydrométrique existe déjà ; l'extension porte sur la modélisation &
moyen \\
Autre compartiment (air, sol) & Les sources et les indicateurs diffèrent, le patron tient &
moyen \\
Couplage à un modèle hydrogéologique distribué & Une interface avec un simulateur physique
externe et l'assimilation de ses sorties & élevé \\
\bottomrule
\end{tabularx}
\end{table}

Ce qui est peu coûteux l'est \emph{parce que} l'infrastructure existe. Le coût élevé de l'année
écoulée a été payé une fois ; il ne se représente pas à chaque extension. C'est la définition
même d'un socle.

\section{Articulation avec le volet des jumeaux numériques}

Un jumeau numérique de la ressource en eau suppose trois choses : une représentation de l'état
courant du système, une capacité à le projeter, et une capacité à expliquer la projection.
L'entrepôt et l'observatoire couvrent l'état, les modèles métier et appris la projection.

La contribution la plus spécifique porte sur la troisième. Un jumeau numérique qui projette sans
expliquer n'est qu'un simulateur ; ce qui en fait un instrument de décision est la capacité à
répondre « pourquoi » et « que se passerait-il si » dans les termes du métier. Les contrefactuels
sous contrainte physique décrits en section~\ref{sec:physcf} sont directement conçus pour cet
usage : ils produisent des scénarios énonçables, du type « il aurait fallu des hivers un
cinquième plus humides et un climat plus frais d'un degré », plutôt que des perturbations
numériques.

\section{État de passation}

Le contrat s'est achevé le 30 juin 2026. Les travaux se sont poursuivis jusqu'au 26 août 2026
dans le seul but de rendre les réalisations exploitables sans leur auteur :

\begin{itemize}
  \item refonte complète de la documentation des deux dépôts, organisée autour d'une
  cartographie des documents et vérifiée par exécution réelle des procédures décrites ;
  \item consignation des pièges susceptibles de produire des résultats faux sans émettre
  d'erreur, développés au chapitre~\ref{ch:entrepot} ;
  \item retrait des éléments sensibles : identifiants, clés d'interface, topologie interne,
  adresses personnelles ; restriction par défaut de l'écoute des services à l'interface locale ;
  \item suppression du code mort et des scripts devenus inopérants ;
  \item vérification à l'échelle réelle des procédures de sauvegarde et de restauration ;
  \item explicitation du couplage entre les deux dépôts et de ses conditions de démarrage.
\end{itemize}

Aucun repreneur n'est désigné à la date de ce rapport. Ce qui précède ne prépare donc pas une
transmission à une personne identifiée : il vise à ce que la reprise reste possible par
quiconque en aura la charge, sans recours à son auteur.

\section{Conclusion}

Ces travaux répondent à une question qui n'était pas celle posée au départ, mais qui la
conditionnait : \emph{comment rendre possible la recherche en intelligence artificielle sur les
eaux souterraines françaises ?}

La réponse tient en une chaîne complète, allant de quatre jeux de données publics et dispersés
jusqu'à des scénarios contrefactuels énonçables dans le langage des hydrogéologues, dont chaque
maillon est documenté, vérifié et réutilisable indépendamment des autres.

Ce qui a été construit dépasse le périmètre des eaux souterraines : l'enchaînement
sources~$\rightarrow$~entrepôt~$\rightarrow$~indicateurs validés~$\rightarrow$~modèles~$\rightarrow$~explication
est transposable aux autres compartiments environnementaux visés par le programme. C'est en ce
sens que ces développements constituent moins un aboutissement qu'un point de départ.
